
%%%%%%%%%%%%%%%%%%%%%%%%%%%%%%%%%%%%%%%%
%           main body
%%%%%%%%%%%%%%%%%%%%%%%%%%%%%%%%%%%%%%%%


%-----------------------------------
%	TITLE PAGE
%-----------------------------------

\title[Probability \& Statistics]{\LARGE \bfseries 第三部分\quad 概率与统计} % The short title appears at the bottom of every slide, the full title is only on the title page
\author[Central Limit Theorem] % Your institution as it will appear on the bottom of every slide, may be shorthand to save space
{\Large \bfseries \S 7. 中心极限定理}
% \author{Singular Value Decomposition} % Your name
% \date{\today} % Date, can be changed to a custom date
\date{}

%------------------------------------------------

\begin{document}


%------------------------------------------------

\begin{frame}

\titlepage % Print the title page as the first slide

\end{frame}

% %------------------------------------------------

% \begin{frame}
% \frametitle{内容提要} % Table of contents slide, comment this block out to remove it
% \tableofcontents % Throughout your presentation, if you choose to use \section{} and \subsection{} commands, these will automatically be printed on this slide as an overview of your presentation
% \end{frame}

%-----------------------------------
%	PRESENTATION SLIDES
%-----------------------------------

%------------------------------------------------

\begin{frame}

{
\large
大数定律说明了独立同分布的随机变量的均值与期望之间的偏差小于任意给定正实数的概率趋于1（依概率收敛）。如果还加上方差有限这个条件，那么我们会有更精细的结果 --- 中心极限定理。
}

\end{frame}

%------------------------------------------------

\begin{frame}

\begin{block}{中心极限定理，Central Limit Theorem}
设$\{X_n\}$是一列独立同分布的随机变量，且有
\[E(X_n) = \mu, \; Var(X_n) = \sigma^2,\]
那么，对于任意实数$x$, 都有
\[\lim\limits_{n\to\infty} P\left( \dfrac{\sum\limits_{i=1}^n X_i - n\mu}{\sqrt{n}\sigma} \le x \right) = \Phi(x) = \dfrac{1}{\sqrt{2\pi}} \int_{-\infty}^x e^{-\frac{t^2}{2}} dt\]
这里的$\Phi(x)$即为标准正态分布的分布函数。
\end{block}

\end{frame}

%------------------------------------------------

\begin{frame}

\begin{block}{中心极限定理的意义}
中心极限定理表明，当$n$充分大时，$n$个具有有限期望与方差的独立同分布的随机变量之和``近似''服从正态分布。虽然一般很难求出$X_1+\cdots+X_n$的分布的确切形式，但当$n$很大时，可得出其``近似分布''：
\[\dfrac{\sum\limits_{i=1}^n X_i - n\mu}{\sqrt{n}\sigma} \; \text{``$\sim$''} \; N(0,1)\]
或等价地
\[\sum\limits_{i=1}^n X_i \; \text{``$\sim$''} \; N(n\mu,n\sigma^2)\]
\end{block}

\end{frame}

%------------------------------------------------

\begin{frame}

将中心极限定理应用到参数为$p$的$n$重伯努利试验，那么有
\vspace{0.5em}
\begin{block}{德莫瓦弗（de Moivre）--拉普拉斯定理}
设随机变量$X$服从二项分布$B(n,p)$, 那么
\[\lim\limits_{n\to\infty} P\left( \dfrac{X - np}{\sqrt{np(1-p)}} \le x \right) = \Phi(x)\]
即当$n$充分大时，$X$近似地服从正态分布$N(np,np(1-p))$.
\end{block}

\end{frame}

%------------------------------------------------

\begin{frame}

\begin{block}{例子}
问题：假设测量一个物理量，每次测量都有微小误差。每次的误差在适当的单位下，平均分布于-1与1之间。也就是说把每次测量的误差设为随机变量$X$，那么$X$服从均匀分布$\mathcal{U}(-1,1)$. 现在我们做$n$次独立的测量，那么测量平均值与真值的差的绝对值小于一个微小正数$\varepsilon$的概率是多少?

\vspace{1em}
\pause

解：设$X_i$为第$i$次测量产生的误差，那么$X_i \sim \mathcal{U}(-1,1)$, 从而有$E(X_i) = 0$, $\sigma(X_i) = \frac{1}{3}$. 那么根据中心极限定理，当$n$充分大时有
\[\dfrac{X_1+\cdots+X_n}{\sqrt{n}\sigma} \text{ 近似服从标准正态分布 } N(0,1)\]
于是，我们要求的概率即为
\[P\left( -\dfrac{\sqrt{n}\varepsilon}{\sigma} < \dfrac{X_1+\cdots+X_n}{\sqrt{n}\sigma} < \dfrac{\sqrt{n}\varepsilon}{\sigma} \right) = 2\Phi\left( \dfrac{\sqrt{n}\varepsilon}{\sigma} \right) - 1\]
\end{block}

\end{frame}

%------------------------------------------------

% \begin{frame}

% \begin{block}{Vandermonde行列式的计算}
% \begin{enumerate}
% % \addtocounter{enumi}{0}
% \item 把第$n-1$行乘以$-a_1$加到第$n$行，再把第$n-2$行乘以$-a_1$加到第$n-1$行。按照这种顺序依次向上，直到用第$2$行乘以$-a_1$加到第$1$行。由行列式性质，$V_n$的值不变。
% \[
% V_n = \det \begin{pmatrix}
% 1 & 1 & \cdots & 1 \\
% 0 & a_2-a_1 & \cdots & a_n-a_1 \\
% \vdots & \vdots & & \vdots \\
% 0 & a_2^{n-2}(a_2-a_1) & \cdots & a_n^{n-2}(a_n-a_1) \\
% \end{pmatrix}
% \]
% \end{enumerate}
% \end{block}

% \end{frame}

%------------------------------------------------
% % closing page
% \begin{frame}
% \HUGE{\centerline{\bfseries {\color{gray} The} {\color{zkblue} End}}}

% \pause

% \vspace{1.2em}

% \Large{\centerline{\bfseries {\color{gray}Thank} \quad {\color{zkblue} You}}}

% \end{frame}

%----------------------------------------------------------------------------------------

\end{document}
